\section{Discussion}\label{sec:discussion}

\subsection{Why our share is lower than Kim's}

Our label is not the reason: Kim's own rules give a share even lower than ours on our data (Section~\ref{sec:robustness}). The choice of which migrations to count matters more. Counting from 30 days before the notice, as Kim does, removes the projects that moved early and raises the share to \AfterWindow\%, and counting every recent shutdown lowers it to \AfterAll\%. The remaining gap may come from the set of models, the mining period and the search method, which differ between the two studies; Kim's paper is a preprint, and some of its details cannot yet be compared. The robust finding is the direction: more projects migrate after the shutdown than before it.

\subsection{Implications}

\paragraph{For developers.} The busiest moment to migrate is the month after the model has stopped answering, so many applications were probably broken for some days unless they had a fallback model. Deprecation notices give months of warning, and projects that migrate during the notice do so when, at the median, \ExactNoticeUsed\% of it has passed. A check in CI that flags model IDs with an announced retirement would turn the notice into a failing build rather than an outage. When they do migrate, projects should check the target: one in five moves to a model that is already on a retirement list.

\paragraph{For providers.} Most projects follow the recommended model line, so the recommendation is heard. Yet the share of late migrations suggests that the notice does not reach many maintainers. Notices delivered where developers look, such as warnings in API responses before the shutdown, may help more than a longer notice period.

\paragraph{For researchers and tool builders.} Only about one migration in six touches a test and very few touch an evaluation, although a new model can change behaviour in ways tests would catch \citep{ma2024prompt}. Tools that generate a behavioural regression check for the old and new model, or that make the switch with a pull request that includes such a check, address a gap that the data shows clearly. The released dataset (see the data availability statement) can serve as a benchmark for such tools.
